% Pursuing Stacks (cotes 134-2 to 134-8): archive cote 134-7, site page 76 (PDF page 77), typescript page 518.
% typescript. Draft compared: ps6.tex:3795-3857.
% AI-transcribed by Claude Opus 5.5 (claude-opus-5-5, adaptive thinking, effort high), Message Batches API, batch ps4 (msgbatch_012Aoc1NuskATPsyBeKSYK7J), 2026-11-12; not yet checked by eye.
\scanpage{134-7/76}
\mspage{518}

4.10.

123.\quad For the last few days  - since I resumed the daily mathematical
ponderings (interrupted more or less for nearly one month),  the wind in
my sails has been rather slack I feel.  It has been nearly six weeks
from now that I started pondering on schematization and on schematic
homotopy types - in the process I got rid of some misplaced expectations,
very well.  Still,  the unpleasant feeling remains of not having really
any hold yet on those would-be schematic homotopy types,  due mainly to
my uncapacity so far of performing (in the model category of semisimpli-%
cial unipotent bundles say) the basic homotopy operations of taking
homotopy fibers and cofibers of maps.  In lack of this,  I am not \deleted{even}
(``morally'') sure\inserted{yet} that there does exist indeed a substantial \deleted{\unsure{the}} reality
of the kind I have been trying to foreshadow.  I am not even wholly clear
yet of how to define the notion of ``weak equivalence'' in the model catego
% NOTE: the right edge of this line and the next is cut off by the scan
ry $\hand{\mathcal{M}}_\bullet$(k) (or in the smaller one $\hand{\mathcal{M}}_1$(k), to play safe, as even the defini
tion of $\hand{\mathcal{M}}_\bullet$(k) is'nt too clear yet) -  there are three ways to define
weak equivalences,  using LH. , or L\hand{$\pi$}. , or the sections functor
from $\hand{\mathcal{M}}_1$(k) to ordinary semisimplicial complexes,  and it is'nt clear
yet whether these are indeed equivalent.  We may of course call ``weak
equivalence'' in $\hand{\mathcal{M}}_1$(k)  a map which \deleted{is transformed to}\inserted{becomes} an isomorphism by
any of these three functors.  But I would'nt say I am wholly confident
yet that there exists \xout{say} a localization Hot$_1$(k) of $\hand{\mathcal{M}}_1$(k)\hand{,} wi th respect
to this notion of weak equivalence or some finer one,  in such a way that
in H\over{e}{o}t$_1$(k) one may introduce the two types of ``fibration'' and ``cofibrati-%
on'' sequences with the usual properties,  and the sections functor
\[ \mathrm{Hot}_1(k) \hand{\longrightarrow} \mathrm{Hot}_1 \]
should respect this structure,  and moreover give rise to functorial
isomorphisms
\[ H_i(L\hand{\pi}.(X_*)) \hand{\simeq} \hand{\pi}_i(X_*(k)) \qquad ; \]
nor even do I feel wholly confident that a theory of this kind can be
developped,  \deleted{using} possibly\inserted{with} different kinds of models for k-homotopy
types from the ones I have been using.  At any rate,  it seems to me \deleted{that}
that the two kind of conditions or features I have just been stating,
are indeed the crucial ones,  plus (I should add)  existence in H\over{e}{o}t$_1$(k)
of an object \deleted{\illegible{1}}
\[ S^2_{\hand{k}} = S(2,\over{\illegible{1}}{k}) \quad , \]
standing for the  2-sphere,  with
\[ \hand{\widetilde{\phantom{LH}}}\!\!\!\!\mathrm{LH}.(S^2_{\hand{k}}) \simeq \underline{Z}[2] , \]

\margin{76}
